
\documentclass{article}
\usepackage{colortbl}
\usepackage{makecell}
\usepackage{multirow}
\usepackage{supertabular}

\begin{document}

\newcounter{utterance}

\twocolumn

{ \footnotesize  \setcounter{utterance}{1}
\setlength{\tabcolsep}{0pt}
\begin{supertabular}{c@{$\;$}|p{.15\linewidth}@{}p{.15\linewidth}p{.15\linewidth}p{.15\linewidth}p{.15\linewidth}p{.15\linewidth}}

    \# & $\;$A & \multicolumn{4}{c}{Game Master} & $\;\:$B\\
    \hline 

    \theutterance \stepcounter{utterance}  

    & & \multicolumn{4}{p{0.6\linewidth}}{\cellcolor[rgb]{0.9,0.9,0.9}{%
	\makecell[{{p{\linewidth}}}]{% 
	  \tt {\tiny [A$\langle$GM]}  
	 Sono il tuo game master, e stai giocando un gioco collaborativo con la seguente griglia come tavolo da gioco:\\ \tt \\ \tt ```\\ \tt  1234567\\ \tt ╔═╤═╤═╤═╗\\ \tt ║◌│◌│◌│◌║ 1\\ \tt ║◌└─┼─┴─╢ 2\\ \tt ║◌◌◌│◌◌◌║ 3\\ \tt ╟──┬┴───╢ 4\\ \tt ║◌◌│P◌◌◌║ 5\\ \tt ║LC│◌◌◌◌║ 6\\ \tt ║◌◌│◌◌◌◌║ 7\\ \tt ╚══╧════╝\\ \tt \\ \tt ```\\ \tt \\ \tt * Il bordo superiore mostra le coordinate x che aumentano verso destra, e il bordo destro le coordinate y che aumentano verso il basso.\\ \tt * I seguenti oggetti sono posizionati casualmente sulla tua griglia: 'C', 'L', 'P'.\\ \tt \\ \tt \\ \tt L'altro giocatore vede una variante del tavolo da gioco, dove gli oggetti sono posizionati in diverse posizioni casuali. Non puoi vedere il tavolo dell'altro giocatore, e loro non possono vedere il tuo.\\ \tt \\ \tt \\ \tt **Obiettivo:**\\ \tt \\ \tt Entrambi i giocatori devono spostare gli oggetti sui rispettivi sfondi in modo che oggetti identici finiscano alle stesse coordinate. Devi comunicare con l'altro giocatore per concordare una configurazione obiettivo comune.\\ \tt \\ \tt \\ \tt **Regole:**\\ \tt \\ \tt * In ogni turno, puoi inviare esattamente uno dei seguenti due comandi:\\ \tt  1. `COMUNICA: <MESSAGE>`: per inviare un messaggio (tutto fino alla prossima interruzione di riga) all'altro giocatore. Lo inoltrerò al tuo partner.\\ \tt 2. `SPOSTA: <OBJECT>, (<X>, <Y>)`: per spostare un oggetto in una nuova posizione, dove <X> è la colonna e <Y> è la riga. Ti farò sapere se la tua mossa è andata a buon fine o no.\\ \tt * Se non rispetti il formato, o invii diversi comandi contemporaneamente, devo penalizzarti.\\ \tt * Se entrambi i giocatori accumulano più di 8 penalità, perdete entrambi il gioco.\\ \tt * È vitale che comunichi con l'altro giocatore riguardo al tuo stato obiettivo! L'*unico* modo in cui puoi trasmettere la tua strategia all'altro giocatore è usando il comando `COMUNICA: <MESSAGE>`!\\ \tt \\ \tt \\ \tt **Spostamento Oggetti**\\ \tt \\ \tt * Puoi spostare oggetti solo in celle entro i limiti della griglia. La cella di destinazione deve essere vuota, cioè deve contenere solo il simbolo '◌'.\\ \tt * Se provi a spostare un oggetto in un posto che non è vuoto, o provi a spostarlo fuori dalla griglia, devo penalizzarti. Hai un altro tentativo.\\ \tt * Prima di fare una mossa, controlla due volte che il posto di destinazione sia vuoto, e non contenga alcuna lettera, cornice, o linea!\\ \tt \\ \tt \\ \tt **Fine del Gioco**\\ \tt \\ \tt Se pensi di aver raggiunto l'obiettivo di allineare tutti gli oggetti, puoi chiedere all'altro giocatore di finire il gioco inviando COMUNICA: finito?. Se l'altro giocatore ti chiede di finire il gioco, e rispondi con COMUNICA: finito!, the game will end.\\ \tt \\ \tt Entrambi i giocatori vincono se il gioco finisce entro 12 round, dove un round è definito come due giocatori che inviano ciascuno un comando valido.\\ \tt \\ \tt **Punteggio:**\\ \tt \\ \tt Più vicini sono gli oggetti identici in entrambi i tavoli da gioco, più punti ottieni. Le penalità riducono i tuoi punti. Riesci a battere il record?\\ \tt \\ \tt \\ \tt Per favore invia un messaggio all'altro giocatore per iniziare il gioco! 
	  } 
	   } 
	   } 
	 & \\ 
 

    \theutterance \stepcounter{utterance}  

    & \multicolumn{4}{p{0.6\linewidth}}{\cellcolor[rgb]{0.8,1,0.9}{%
	\makecell[{{p{\linewidth}}}]{% 
	  \tt {\tiny [A$\rangle$GM]}  
	 COMUNICA: Ciao! Sto cercando di coordinare le posizioni degli oggetti. Possiamo concordare un piano per spostare 'C', 'L' e 'P' verso posizioni specifiche? Per favore, condividi le tue idee o suggerisci un punto di partenza. 😊 
	  } 
	   } 
	   } 
	 & & \\ 
 

    \theutterance \stepcounter{utterance}  

    & & & \multicolumn{2}{p{0.3\linewidth}}{\cellcolor[rgb]{0.95,0.95,0.95}{%
	\makecell[{{p{\linewidth}}}]{% 
	  \tt {\tiny [GM$|$GM]}  
	 COMUNICA: Ciao! Sto cercando di coordinare le posizioni degli oggetti. Possiamo concordare un piano per spostare 'C', 'L' e 'P' verso posizioni specifiche? Per favore, condividi le tue idee o suggerisci un punto di partenza. 😊 
	  } 
	   } 
	   } 
	 & & \\ 
 

    \theutterance \stepcounter{utterance}  

    & & \multicolumn{4}{p{0.6\linewidth}}{\cellcolor[rgb]{0.9,0.9,0.9}{%
	\makecell[{{p{\linewidth}}}]{% 
	  \tt {\tiny [GM$\rangle$B]}  
	 Sono il tuo game master, e stai giocando un gioco collaborativo con la seguente griglia come tavolo da gioco:\\ \tt \\ \tt ```\\ \tt  1234567\\ \tt ╔═╤═╤═╤═╗\\ \tt ║◌│◌│L│◌║ 1\\ \tt ║◌└─┼─┴─╢ 2\\ \tt ║◌◌◌│◌◌C║ 3\\ \tt ╟──┬┴───╢ 4\\ \tt ║◌◌│◌◌◌◌║ 5\\ \tt ║◌◌│◌P◌◌║ 6\\ \tt ║◌◌│◌◌◌◌║ 7\\ \tt ╚══╧════╝\\ \tt \\ \tt ```\\ \tt \\ \tt * Il bordo superiore mostra le coordinate x che aumentano verso destra, e il bordo destro le coordinate y che aumentano verso il basso.\\ \tt * I seguenti oggetti sono posizionati casualmente sulla tua griglia: 'C', 'L', 'P'.\\ \tt \\ \tt \\ \tt L'altro giocatore vede una variante del tavolo da gioco, dove gli oggetti sono posizionati in diverse posizioni casuali. Non puoi vedere il tavolo dell'altro giocatore, e loro non possono vedere il tuo.\\ \tt \\ \tt \\ \tt **Obiettivo:**\\ \tt \\ \tt Entrambi i giocatori devono spostare gli oggetti sui rispettivi sfondi in modo che oggetti identici finiscano alle stesse coordinate. Devi comunicare con l'altro giocatore per concordare una configurazione obiettivo comune.\\ \tt \\ \tt \\ \tt **Regole:**\\ \tt \\ \tt * In ogni turno, puoi inviare esattamente uno dei seguenti due comandi:\\ \tt  1. `COMUNICA: <MESSAGE>`: per inviare un messaggio (tutto fino alla prossima interruzione di riga) all'altro giocatore. Lo inoltrerò al tuo partner.\\ \tt 2. `SPOSTA: <OBJECT>, (<X>, <Y>)`: per spostare un oggetto in una nuova posizione, dove <X> è la colonna e <Y> è la riga. Ti farò sapere se la tua mossa è andata a buon fine o no.\\ \tt * Se non rispetti il formato, o invii diversi comandi contemporaneamente, devo penalizzarti.\\ \tt * Se entrambi i giocatori accumulano più di 8 penalità, perdete entrambi il gioco.\\ \tt * È vitale che comunichi con l'altro giocatore riguardo al tuo stato obiettivo! L'*unico* modo in cui puoi trasmettere la tua strategia all'altro giocatore è usando il comando `COMUNICA: <MESSAGE>`!\\ \tt \\ \tt \\ \tt **Spostamento Oggetti**\\ \tt \\ \tt * Puoi spostare oggetti solo in celle entro i limiti della griglia. La cella di destinazione deve essere vuota, cioè deve contenere solo il simbolo '◌'.\\ \tt * Se provi a spostare un oggetto in un posto che non è vuoto, o provi a spostarlo fuori dalla griglia, devo penalizzarti. Hai un altro tentativo.\\ \tt * Prima di fare una mossa, controlla due volte che il posto di destinazione sia vuoto, e non contenga alcuna lettera, cornice, o linea!\\ \tt \\ \tt \\ \tt **Fine del Gioco**\\ \tt \\ \tt Se pensi di aver raggiunto l'obiettivo di allineare tutti gli oggetti, puoi chiedere all'altro giocatore di finire il gioco inviando COMUNICA: finito?. Se l'altro giocatore ti chiede di finire il gioco, e rispondi con COMUNICA: finito!, the game will end.\\ \tt \\ \tt Entrambi i giocatori vincono se il gioco finisce entro 12 round, dove un round è definito come due giocatori che inviano ciascuno un comando valido.\\ \tt \\ \tt **Punteggio:**\\ \tt \\ \tt Più vicini sono gli oggetti identici in entrambi i tavoli da gioco, più punti ottieni. Le penalità riducono i tuoi punti. Riesci a battere il record?\\ \tt \\ \tt L'altro giocatore ha iniziato il gioco inviando questo messaggio:\\ \tt \\ \tt "Ciao! Sto cercando di coordinare le posizioni degli oggetti. Possiamo concordare un piano per spostare 'C', 'L' e 'P' verso posizioni specifiche? Per favore, condividi le tue idee o suggerisci un punto di partenza. 😊"\\ \tt \\ \tt Qual è il tuo primo comando? 
	  } 
	   } 
	   } 
	 & \\ 
 

    \theutterance \stepcounter{utterance}  

    & & & \multicolumn{4}{p{0.6\linewidth}}{\cellcolor[rgb]{1,0.85,0.72}{%
	\makecell[{{p{\linewidth}}}]{% 
	  \tt {\tiny [GM$\langle$B]}  
	 COMUNICA: Ciao! Sto anche io cercando di coordinare le posizioni. Possiamo iniziare concordando un punto di riferimento comune per ogni oggetto. Ad esempio, possiamo decidere di spostare 'C' alla posizione (4,3), 'L' a (2,1) e 'P' a (4,6). Se queste posizioni non sono disponibili, possiamo adattarle. Quali sono le posizioni iniziali degli oggetti sul tuo tavolo? 😊 
	  } 
	   } 
	   } 
	  \\ 
 

    \theutterance \stepcounter{utterance}  

    & & & \multicolumn{2}{p{0.3\linewidth}}{\cellcolor[rgb]{0.95,0.95,0.95}{%
	\makecell[{{p{\linewidth}}}]{% 
	  \tt {\tiny [GM$|$GM]}  
	 COMUNICA: Ciao! Sto anche io cercando di coordinare le posizioni. Possiamo iniziare concordando un punto di riferimento comune per ogni oggetto. Ad esempio, possiamo decidere di spostare 'C' alla posizione (4,3), 'L' a (2,1) e 'P' a (4,6). Se queste posizioni non sono disponibili, possiamo adattarle. Quali sono le posizioni iniziali degli oggetti sul tuo tavolo? 😊 
	  } 
	   } 
	   } 
	 & & \\ 
 

    \theutterance \stepcounter{utterance}  

    & & \multicolumn{4}{p{0.6\linewidth}}{\cellcolor[rgb]{0.9,0.9,0.9}{%
	\makecell[{{p{\linewidth}}}]{% 
	  \tt {\tiny [A$\langle$GM]}  
	 Il tuo messaggio è stato inoltrato all'altro giocatore.\\ \tt \\ \tt \\ \tt Stai attualmente giocando il round 1 di massimo 12.\\ \tt Avete accumulato collettivamente 0 di 8 penalità.\\ \tt L'altro giocatore ha inviato questo messaggio:\\ \tt \\ \tt \\ \tt "Ciao! Sto anche io cercando di coordinare le posizioni. Possiamo iniziare concordando un punto di riferimento comune per ogni oggetto. Ad esempio, possiamo decidere di spostare 'C' alla posizione (4,3), 'L' a (2,1) e 'P' a (4,6). Se queste posizioni non sono disponibili, possiamo adattarle. Quali sono le posizioni iniziali degli oggetti sul tuo tavolo? 😊"\\ \tt \\ \tt \\ \tt Qual è il tuo prossimo comando? 
	  } 
	   } 
	   } 
	 & \\ 
 

    \theutterance \stepcounter{utterance}  

    & \multicolumn{4}{p{0.6\linewidth}}{\cellcolor[rgb]{0.8,1,0.9}{%
	\makecell[{{p{\linewidth}}}]{% 
	  \tt {\tiny [A$\rangle$GM]}  
	 COMUNICA: Grazie per la proposta! Sul mio tavolo, le posizioni iniziali sono:  \\ \tt - 'C' a (3,2)  \\ \tt - 'L' a (2,3)  \\ \tt - 'P' a (5,5)  \\ \tt \\ \tt Le tue posizioni proposte per 'C' (4,3) e 'L' (2,1) sono valide, ma 'P' a (4,6) è fuori dalla griglia (la riga 6 esiste, ma la colonna 4 non è valida per la riga 6, poiché la griglia ha solo 7 colonne e la riga 6 è limitata a colonne 1-7). Possiamo adattare 'P' a (4,5) invece?  \\ \tt \\ \tt Se concordiamo su:  \\ \tt - 'C' → (4,3)  \\ \tt - 'L' → (2,1)  \\ \tt - 'P' → (4,5)  \\ \tt \\ \tt Posso iniziare a spostare gli oggetti. Vuoi che proceda con questi obiettivi? 
	  } 
	   } 
	   } 
	 & & \\ 
 

    \theutterance \stepcounter{utterance}  

    & & & \multicolumn{2}{p{0.3\linewidth}}{\cellcolor[rgb]{0.95,0.95,0.95}{%
	\makecell[{{p{\linewidth}}}]{% 
	  \tt {\tiny [GM$|$GM]}  
	 COMUNICA: Grazie per la proposta! Sul mio tavolo, le posizioni iniziali sono:  \\ \tt - 'C' a (3,2)  \\ \tt - 'L' a (2,3)  \\ \tt - 'P' a (5,5)  \\ \tt \\ \tt Le tue posizioni proposte per 'C' (4,3) e 'L' (2,1) sono valide, ma 'P' a (4,6) è fuori dalla griglia (la riga 6 esiste, ma la colonna 4 non è valida per la riga 6, poiché la griglia ha solo 7 colonne e la riga 6 è limitata a colonne 1-7). Possiamo adattare 'P' a (4,5) invece?  \\ \tt \\ \tt Se concordiamo su:  \\ \tt - 'C' → (4,3)  \\ \tt - 'L' → (2,1)  \\ \tt - 'P' → (4,5)  \\ \tt \\ \tt Posso iniziare a spostare gli oggetti. Vuoi che proceda con questi obiettivi? 
	  } 
	   } 
	   } 
	 & & \\ 
 

    \theutterance \stepcounter{utterance}  

    & & & \multicolumn{2}{p{0.3\linewidth}}{\cellcolor[rgb]{0.95,0.95,0.95}{%
	\makecell[{{p{\linewidth}}}]{% 
	  \tt {\tiny [GM$|$GM]}  
	 Invalid format: tail is not empty: '\\ \tt - 'C' a (3,2)  \\ \tt - 'L' a (2,3)  \\ \tt - 'P' a (5,5)  \\ \tt \\ \tt Le tue posizioni proposte per 'C' (4,3) e 'L' (2,1) sono valide, ma 'P' a (4,6) è fuori dalla griglia (la riga 6 esiste, ma la colonna 4 non è valida per la riga 6, poiché la griglia ha solo 7 colonne e la riga 6 è limitata a colonne 1-7). Possiamo adattare 'P' a (4,5) invece?  \\ \tt \\ \tt Se concordiamo su:  \\ \tt - 'C' → (4,3)  \\ \tt - 'L' → (2,1)  \\ \tt - 'P' → (4,5)  \\ \tt \\ \tt Posso iniziare a spostare gli oggetti. Vuoi che proceda con questi obiettivi? 
	  } 
	   } 
	   } 
	 & & \\ 
 

    \theutterance \stepcounter{utterance}  

    & & \multicolumn{4}{p{0.6\linewidth}}{\cellcolor[rgb]{0.9,0.9,0.9}{%
	\makecell[{{p{\linewidth}}}]{% 
	  \tt {\tiny [A$\langle$GM]}  
	 Penalità: Il tuo messaggio non deve contenere nulla dopo il comando!\\ \tt \\ \tt Assicurati che la tua risposta contenga solo `COMUNICA: <MESSAGE>` o `SPOSTA: <OBJECT>, (<X>, <Y>)`, e nient'altro!\\ \tt \\ \tt Avete accumulato collettivamente 1 di 8 penalità.\\ \tt Riprova! 
	  } 
	   } 
	   } 
	 & \\ 
 

    \theutterance \stepcounter{utterance}  

    & \multicolumn{4}{p{0.6\linewidth}}{\cellcolor[rgb]{0.8,1,0.9}{%
	\makecell[{{p{\linewidth}}}]{% 
	  \tt {\tiny [A$\rangle$GM]}  
	 COMUNICA: Accordo sulle posizioni: 'C' a (4,3), 'L' a (2,1), 'P' a (4,5). Procedo con i movimenti. 
	  } 
	   } 
	   } 
	 & & \\ 
 

    \theutterance \stepcounter{utterance}  

    & & & \multicolumn{2}{p{0.3\linewidth}}{\cellcolor[rgb]{0.95,0.95,0.95}{%
	\makecell[{{p{\linewidth}}}]{% 
	  \tt {\tiny [GM$|$GM]}  
	 COMUNICA: Accordo sulle posizioni: 'C' a (4,3), 'L' a (2,1), 'P' a (4,5). Procedo con i movimenti. 
	  } 
	   } 
	   } 
	 & & \\ 
 

    \theutterance \stepcounter{utterance}  

    & & \multicolumn{4}{p{0.6\linewidth}}{\cellcolor[rgb]{0.9,0.9,0.9}{%
	\makecell[{{p{\linewidth}}}]{% 
	  \tt {\tiny [GM$\rangle$B]}  
	 Il tuo messaggio è stato inoltrato all'altro giocatore.\\ \tt \\ \tt \\ \tt Stai attualmente giocando il round 2 di massimo 12.\\ \tt Avete accumulato collettivamente 1 di 8 penalità.\\ \tt L'altro giocatore ha inviato questo messaggio:\\ \tt \\ \tt \\ \tt "Accordo sulle posizioni: 'C' a (4,3), 'L' a (2,1), 'P' a (4,5). Procedo con i movimenti."\\ \tt \\ \tt \\ \tt Qual è il tuo prossimo comando? 
	  } 
	   } 
	   } 
	 & \\ 
 

    \theutterance \stepcounter{utterance}  

    & & & \multicolumn{4}{p{0.6\linewidth}}{\cellcolor[rgb]{1,0.85,0.72}{%
	\makecell[{{p{\linewidth}}}]{% 
	  \tt {\tiny [GM$\langle$B]}  
	 COMUNICA: Accordo sulle posizioni: 'C' a (4,3), 'L' a (2,1), 'P' a (4,5). Procedo con i movimenti. Verifica se queste coordinate sono disponibili sul tuo tavolo e conferma. 😊 
	  } 
	   } 
	   } 
	  \\ 
 

    \theutterance \stepcounter{utterance}  

    & & & \multicolumn{2}{p{0.3\linewidth}}{\cellcolor[rgb]{0.95,0.95,0.95}{%
	\makecell[{{p{\linewidth}}}]{% 
	  \tt {\tiny [GM$|$GM]}  
	 COMUNICA: Accordo sulle posizioni: 'C' a (4,3), 'L' a (2,1), 'P' a (4,5). Procedo con i movimenti. Verifica se queste coordinate sono disponibili sul tuo tavolo e conferma. 😊 
	  } 
	   } 
	   } 
	 & & \\ 
 

    \theutterance \stepcounter{utterance}  

    & & \multicolumn{4}{p{0.6\linewidth}}{\cellcolor[rgb]{0.9,0.9,0.9}{%
	\makecell[{{p{\linewidth}}}]{% 
	  \tt {\tiny [A$\langle$GM]}  
	 Il tuo messaggio è stato inoltrato all'altro giocatore.\\ \tt \\ \tt \\ \tt Stai attualmente giocando il round 2 di massimo 12.\\ \tt Avete accumulato collettivamente 1 di 8 penalità.\\ \tt L'altro giocatore ha inviato questo messaggio:\\ \tt \\ \tt \\ \tt "Accordo sulle posizioni: 'C' a (4,3), 'L' a (2,1), 'P' a (4,5). Procedo con i movimenti. Verifica se queste coordinate sono disponibili sul tuo tavolo e conferma. 😊"\\ \tt \\ \tt \\ \tt Qual è il tuo prossimo comando? 
	  } 
	   } 
	   } 
	 & \\ 
 

    \theutterance \stepcounter{utterance}  

    & \multicolumn{4}{p{0.6\linewidth}}{\cellcolor[rgb]{0.8,1,0.9}{%
	\makecell[{{p{\linewidth}}}]{% 
	  \tt {\tiny [A$\rangle$GM]}  
	 SPOSTA: C, (4,3) 
	  } 
	   } 
	   } 
	 & & \\ 
 

    \theutterance \stepcounter{utterance}  

    & & & \multicolumn{2}{p{0.3\linewidth}}{\cellcolor[rgb]{0.95,0.95,0.95}{%
	\makecell[{{p{\linewidth}}}]{% 
	  \tt {\tiny [GM$|$GM]}  
	 SPOSTA: C, (4,3) 
	  } 
	   } 
	   } 
	 & & \\ 
 

    \theutterance \stepcounter{utterance}  

    & & & \multicolumn{2}{p{0.3\linewidth}}{\cellcolor[rgb]{0.95,0.95,0.95}{%
	\makecell[{{p{\linewidth}}}]{% 
	  \tt {\tiny [GM$|$GM]}  
	 Penalità: (4,3) non è vuoto, ma contiene '│'.\\ \tt \\ \tt Avete accumulato collettivamente 2 di 8 penalità.\\ \tt Riprova! 
	  } 
	   } 
	   } 
	 & & \\ 
 

    \theutterance \stepcounter{utterance}  

    & & \multicolumn{4}{p{0.6\linewidth}}{\cellcolor[rgb]{0.9,0.9,0.9}{%
	\makecell[{{p{\linewidth}}}]{% 
	  \tt {\tiny [A$\langle$GM]}  
	 Penalità: (4,3) non è vuoto, ma contiene '│'.\\ \tt \\ \tt Avete accumulato collettivamente 2 di 8 penalità.\\ \tt Riprova! 
	  } 
	   } 
	   } 
	 & \\ 
 

    \theutterance \stepcounter{utterance}  

    & \multicolumn{4}{p{0.6\linewidth}}{\cellcolor[rgb]{0.8,1,0.9}{%
	\makecell[{{p{\linewidth}}}]{% 
	  \tt {\tiny [A$\rangle$GM]}  
	 SPOSTA: C, (3,2) 
	  } 
	   } 
	   } 
	 & & \\ 
 

    \theutterance \stepcounter{utterance}  

    & & & \multicolumn{2}{p{0.3\linewidth}}{\cellcolor[rgb]{0.95,0.95,0.95}{%
	\makecell[{{p{\linewidth}}}]{% 
	  \tt {\tiny [GM$|$GM]}  
	 SPOSTA: C, (3,2) 
	  } 
	   } 
	   } 
	 & & \\ 
 

    \theutterance \stepcounter{utterance}  

    & & & \multicolumn{2}{p{0.3\linewidth}}{\cellcolor[rgb]{0.95,0.95,0.95}{%
	\makecell[{{p{\linewidth}}}]{% 
	  \tt {\tiny [GM$|$GM]}  
	 Penalità: (3,2) non è vuoto, ma contiene '─'.\\ \tt \\ \tt Avete accumulato collettivamente 3 di 8 penalità.\\ \tt Riprova! 
	  } 
	   } 
	   } 
	 & & \\ 
 

    \theutterance \stepcounter{utterance}  

    & & \multicolumn{4}{p{0.6\linewidth}}{\cellcolor[rgb]{0.9,0.9,0.9}{%
	\makecell[{{p{\linewidth}}}]{% 
	  \tt {\tiny [A$\langle$GM]}  
	 Penalità: (3,2) non è vuoto, ma contiene '─'.\\ \tt \\ \tt Avete accumulato collettivamente 3 di 8 penalità.\\ \tt Riprova! 
	  } 
	   } 
	   } 
	 & \\ 
 

    \theutterance \stepcounter{utterance}  

    & \multicolumn{4}{p{0.6\linewidth}}{\cellcolor[rgb]{0.8,1,0.9}{%
	\makecell[{{p{\linewidth}}}]{% 
	  \tt {\tiny [A$\rangle$GM]}  
	 SPOSTA: C, (4,4) 
	  } 
	   } 
	   } 
	 & & \\ 
 

    \theutterance \stepcounter{utterance}  

    & & & \multicolumn{2}{p{0.3\linewidth}}{\cellcolor[rgb]{0.95,0.95,0.95}{%
	\makecell[{{p{\linewidth}}}]{% 
	  \tt {\tiny [GM$|$GM]}  
	 SPOSTA: C, (4,4) 
	  } 
	   } 
	   } 
	 & & \\ 
 

    \theutterance \stepcounter{utterance}  

    & & & \multicolumn{2}{p{0.3\linewidth}}{\cellcolor[rgb]{0.95,0.95,0.95}{%
	\makecell[{{p{\linewidth}}}]{% 
	  \tt {\tiny [GM$|$GM]}  
	 Penalità: (4,4) non è vuoto, ma contiene '┴'.\\ \tt \\ \tt Avete accumulato collettivamente 4 di 8 penalità.\\ \tt Riprova! 
	  } 
	   } 
	   } 
	 & & \\ 
 

    \theutterance \stepcounter{utterance}  

    & & \multicolumn{4}{p{0.6\linewidth}}{\cellcolor[rgb]{0.9,0.9,0.9}{%
	\makecell[{{p{\linewidth}}}]{% 
	  \tt {\tiny [A$\langle$GM]}  
	 Penalità: (4,4) non è vuoto, ma contiene '┴'.\\ \tt \\ \tt Avete accumulato collettivamente 4 di 8 penalità.\\ \tt Riprova! 
	  } 
	   } 
	   } 
	 & \\ 
 

    \theutterance \stepcounter{utterance}  

    & \multicolumn{4}{p{0.6\linewidth}}{\cellcolor[rgb]{0.8,1,0.9}{%
	\makecell[{{p{\linewidth}}}]{% 
	  \tt {\tiny [A$\rangle$GM]}  
	 SPOSTA: C, (5,3) 
	  } 
	   } 
	   } 
	 & & \\ 
 

    \theutterance \stepcounter{utterance}  

    & & & \multicolumn{2}{p{0.3\linewidth}}{\cellcolor[rgb]{0.95,0.95,0.95}{%
	\makecell[{{p{\linewidth}}}]{% 
	  \tt {\tiny [GM$|$GM]}  
	 SPOSTA: C, (5,3) 
	  } 
	   } 
	   } 
	 & & \\ 
 

    \theutterance \stepcounter{utterance}  

    & & & \multicolumn{2}{p{0.3\linewidth}}{\cellcolor[rgb]{0.95,0.95,0.95}{%
	\makecell[{{p{\linewidth}}}]{% 
	  \tt {\tiny [GM$|$GM]}  
	 Spostato 'C' in (5,3) con successo. La tua griglia aggiornata appare così:\\ \tt ```\\ \tt  1234567\\ \tt ╔═╤═╤═╤═╗\\ \tt ║◌│◌│◌│◌║ 1\\ \tt ║◌└─┼─┴─╢ 2\\ \tt ║◌◌◌│C◌◌║ 3\\ \tt ╟──┬┴───╢ 4\\ \tt ║◌◌│P◌◌◌║ 5\\ \tt ║L◌│◌◌◌◌║ 6\\ \tt ║◌◌│◌◌◌◌║ 7\\ \tt ╚══╧════╝\\ \tt ```\\ \tt  
	  } 
	   } 
	   } 
	 & & \\ 
 

    \theutterance \stepcounter{utterance}  

    & & \multicolumn{4}{p{0.6\linewidth}}{\cellcolor[rgb]{0.9,0.9,0.9}{%
	\makecell[{{p{\linewidth}}}]{% 
	  \tt {\tiny [GM$\rangle$B]}  
	 Il tuo messaggio è stato inoltrato all'altro giocatore.\\ \tt \\ \tt \\ \tt Stai attualmente giocando il round 3 di massimo 12.\\ \tt Avete accumulato collettivamente 4 di 8 penalità.\\ \tt L'altro giocatore ha spostato un oggetto sulla sua griglia.\\ \tt \\ \tt \\ \tt Qual è il tuo prossimo comando? 
	  } 
	   } 
	   } 
	 & \\ 
 

    \theutterance \stepcounter{utterance}  

    & & & \multicolumn{4}{p{0.6\linewidth}}{\cellcolor[rgb]{1,0.85,0.72}{%
	\makecell[{{p{\linewidth}}}]{% 
	  \tt {\tiny [GM$\langle$B]}  
	 COMUNICA: Grazie per il movimento! Verifica se l'oggetto spostato è in una posizione che si avvicina al nostro obiettivo concordato. Se necessario, possiamo adattare le coordinate per 'C', 'L' o 'P' in base al tuo nuovo stato. 😊 
	  } 
	   } 
	   } 
	  \\ 
 

    \theutterance \stepcounter{utterance}  

    & & & \multicolumn{2}{p{0.3\linewidth}}{\cellcolor[rgb]{0.95,0.95,0.95}{%
	\makecell[{{p{\linewidth}}}]{% 
	  \tt {\tiny [GM$|$GM]}  
	 COMUNICA: Grazie per il movimento! Verifica se l'oggetto spostato è in una posizione che si avvicina al nostro obiettivo concordato. Se necessario, possiamo adattare le coordinate per 'C', 'L' o 'P' in base al tuo nuovo stato. 😊 
	  } 
	   } 
	   } 
	 & & \\ 
 

    \theutterance \stepcounter{utterance}  

    & & \multicolumn{4}{p{0.6\linewidth}}{\cellcolor[rgb]{0.9,0.9,0.9}{%
	\makecell[{{p{\linewidth}}}]{% 
	  \tt {\tiny [A$\langle$GM]}  
	 Spostato 'C' in (5,3) con successo. La tua griglia aggiornata appare così:\\ \tt ```\\ \tt  1234567\\ \tt ╔═╤═╤═╤═╗\\ \tt ║◌│◌│◌│◌║ 1\\ \tt ║◌└─┼─┴─╢ 2\\ \tt ║◌◌◌│C◌◌║ 3\\ \tt ╟──┬┴───╢ 4\\ \tt ║◌◌│P◌◌◌║ 5\\ \tt ║L◌│◌◌◌◌║ 6\\ \tt ║◌◌│◌◌◌◌║ 7\\ \tt ╚══╧════╝\\ \tt ```\\ \tt Stai attualmente giocando il round 3 di massimo 12.\\ \tt Avete accumulato collettivamente 4 di 8 penalità.\\ \tt L'altro giocatore ha inviato questo messaggio:\\ \tt \\ \tt \\ \tt "Grazie per il movimento! Verifica se l'oggetto spostato è in una posizione che si avvicina al nostro obiettivo concordato. Se necessario, possiamo adattare le coordinate per 'C', 'L' o 'P' in base al tuo nuovo stato. 😊"\\ \tt \\ \tt \\ \tt Qual è il tuo prossimo comando? 
	  } 
	   } 
	   } 
	 & \\ 
 

    \theutterance \stepcounter{utterance}  

    & \multicolumn{4}{p{0.6\linewidth}}{\cellcolor[rgb]{0.8,1,0.9}{%
	\makecell[{{p{\linewidth}}}]{% 
	  \tt {\tiny [A$\rangle$GM]}  
	 SPOSTA: L, (2,1) 
	  } 
	   } 
	   } 
	 & & \\ 
 

    \theutterance \stepcounter{utterance}  

    & & & \multicolumn{2}{p{0.3\linewidth}}{\cellcolor[rgb]{0.95,0.95,0.95}{%
	\makecell[{{p{\linewidth}}}]{% 
	  \tt {\tiny [GM$|$GM]}  
	 SPOSTA: L, (2,1) 
	  } 
	   } 
	   } 
	 & & \\ 
 

    \theutterance \stepcounter{utterance}  

    & & & \multicolumn{2}{p{0.3\linewidth}}{\cellcolor[rgb]{0.95,0.95,0.95}{%
	\makecell[{{p{\linewidth}}}]{% 
	  \tt {\tiny [GM$|$GM]}  
	 Penalità: (2,1) non è vuoto, ma contiene '│'.\\ \tt \\ \tt Avete accumulato collettivamente 5 di 8 penalità.\\ \tt Riprova! 
	  } 
	   } 
	   } 
	 & & \\ 
 

    \theutterance \stepcounter{utterance}  

    & & \multicolumn{4}{p{0.6\linewidth}}{\cellcolor[rgb]{0.9,0.9,0.9}{%
	\makecell[{{p{\linewidth}}}]{% 
	  \tt {\tiny [A$\langle$GM]}  
	 Penalità: (2,1) non è vuoto, ma contiene '│'.\\ \tt \\ \tt Avete accumulato collettivamente 5 di 8 penalità.\\ \tt Riprova! 
	  } 
	   } 
	   } 
	 & \\ 
 

    \theutterance \stepcounter{utterance}  

    & \multicolumn{4}{p{0.6\linewidth}}{\cellcolor[rgb]{0.8,1,0.9}{%
	\makecell[{{p{\linewidth}}}]{% 
	  \tt {\tiny [A$\rangle$GM]}  
	 SPOSTA: L, (1,1) 
	  } 
	   } 
	   } 
	 & & \\ 
 

    \theutterance \stepcounter{utterance}  

    & & & \multicolumn{2}{p{0.3\linewidth}}{\cellcolor[rgb]{0.95,0.95,0.95}{%
	\makecell[{{p{\linewidth}}}]{% 
	  \tt {\tiny [GM$|$GM]}  
	 SPOSTA: L, (1,1) 
	  } 
	   } 
	   } 
	 & & \\ 
 

    \theutterance \stepcounter{utterance}  

    & & & \multicolumn{2}{p{0.3\linewidth}}{\cellcolor[rgb]{0.95,0.95,0.95}{%
	\makecell[{{p{\linewidth}}}]{% 
	  \tt {\tiny [GM$|$GM]}  
	 Spostato 'L' in (1,1) con successo. La tua griglia aggiornata appare così:\\ \tt ```\\ \tt  1234567\\ \tt ╔═╤═╤═╤═╗\\ \tt ║L│◌│◌│◌║ 1\\ \tt ║◌└─┼─┴─╢ 2\\ \tt ║◌◌◌│C◌◌║ 3\\ \tt ╟──┬┴───╢ 4\\ \tt ║◌◌│P◌◌◌║ 5\\ \tt ║◌◌│◌◌◌◌║ 6\\ \tt ║◌◌│◌◌◌◌║ 7\\ \tt ╚══╧════╝\\ \tt ```\\ \tt  
	  } 
	   } 
	   } 
	 & & \\ 
 

    \theutterance \stepcounter{utterance}  

    & & \multicolumn{4}{p{0.6\linewidth}}{\cellcolor[rgb]{0.9,0.9,0.9}{%
	\makecell[{{p{\linewidth}}}]{% 
	  \tt {\tiny [GM$\rangle$B]}  
	 Il tuo messaggio è stato inoltrato all'altro giocatore.\\ \tt \\ \tt \\ \tt Stai attualmente giocando il round 4 di massimo 12.\\ \tt Avete accumulato collettivamente 5 di 8 penalità.\\ \tt L'altro giocatore ha spostato un oggetto sulla sua griglia.\\ \tt \\ \tt \\ \tt Qual è il tuo prossimo comando? 
	  } 
	   } 
	   } 
	 & \\ 
 

    \theutterance \stepcounter{utterance}  

    & & & \multicolumn{4}{p{0.6\linewidth}}{\cellcolor[rgb]{1,0.85,0.72}{%
	\makecell[{{p{\linewidth}}}]{% 
	  \tt {\tiny [GM$\langle$B]}  
	 COMUNICA: Grazie per il movimento! Conferma se l'oggetto spostato è in una posizione che si avvicina al nostro obiettivo (C: (4,3), L: (2,1), P: (4,5)). Se non è così, possiamo modificare le coordinate per ottimizzare il piano. 😊 
	  } 
	   } 
	   } 
	  \\ 
 

    \theutterance \stepcounter{utterance}  

    & & & \multicolumn{2}{p{0.3\linewidth}}{\cellcolor[rgb]{0.95,0.95,0.95}{%
	\makecell[{{p{\linewidth}}}]{% 
	  \tt {\tiny [GM$|$GM]}  
	 COMUNICA: Grazie per il movimento! Conferma se l'oggetto spostato è in una posizione che si avvicina al nostro obiettivo (C: (4,3), L: (2,1), P: (4,5)). Se non è così, possiamo modificare le coordinate per ottimizzare il piano. 😊 
	  } 
	   } 
	   } 
	 & & \\ 
 

    \theutterance \stepcounter{utterance}  

    & & \multicolumn{4}{p{0.6\linewidth}}{\cellcolor[rgb]{0.9,0.9,0.9}{%
	\makecell[{{p{\linewidth}}}]{% 
	  \tt {\tiny [A$\langle$GM]}  
	 Spostato 'L' in (1,1) con successo. La tua griglia aggiornata appare così:\\ \tt ```\\ \tt  1234567\\ \tt ╔═╤═╤═╤═╗\\ \tt ║L│◌│◌│◌║ 1\\ \tt ║◌└─┼─┴─╢ 2\\ \tt ║◌◌◌│C◌◌║ 3\\ \tt ╟──┬┴───╢ 4\\ \tt ║◌◌│P◌◌◌║ 5\\ \tt ║◌◌│◌◌◌◌║ 6\\ \tt ║◌◌│◌◌◌◌║ 7\\ \tt ╚══╧════╝\\ \tt ```\\ \tt Stai attualmente giocando il round 4 di massimo 12.\\ \tt Avete accumulato collettivamente 5 di 8 penalità.\\ \tt L'altro giocatore ha inviato questo messaggio:\\ \tt \\ \tt \\ \tt "Grazie per il movimento! Conferma se l'oggetto spostato è in una posizione che si avvicina al nostro obiettivo (C: (4,3), L: (2,1), P: (4,5)). Se non è così, possiamo modificare le coordinate per ottimizzare il piano. 😊"\\ \tt \\ \tt \\ \tt Qual è il tuo prossimo comando? 
	  } 
	   } 
	   } 
	 & \\ 
 

    \theutterance \stepcounter{utterance}  

    & \multicolumn{4}{p{0.6\linewidth}}{\cellcolor[rgb]{0.8,1,0.9}{%
	\makecell[{{p{\linewidth}}}]{% 
	  \tt {\tiny [A$\rangle$GM]}  
	 SPOSTA: P, (4,5) 
	  } 
	   } 
	   } 
	 & & \\ 
 

    \theutterance \stepcounter{utterance}  

    & & & \multicolumn{2}{p{0.3\linewidth}}{\cellcolor[rgb]{0.95,0.95,0.95}{%
	\makecell[{{p{\linewidth}}}]{% 
	  \tt {\tiny [GM$|$GM]}  
	 SPOSTA: P, (4,5) 
	  } 
	   } 
	   } 
	 & & \\ 
 

    \theutterance \stepcounter{utterance}  

    & & & \multicolumn{2}{p{0.3\linewidth}}{\cellcolor[rgb]{0.95,0.95,0.95}{%
	\makecell[{{p{\linewidth}}}]{% 
	  \tt {\tiny [GM$|$GM]}  
	 Penalità: (4,5) non è vuoto, ma contiene 'P'.\\ \tt \\ \tt Avete accumulato collettivamente 6 di 8 penalità.\\ \tt Riprova! 
	  } 
	   } 
	   } 
	 & & \\ 
 

    \theutterance \stepcounter{utterance}  

    & & \multicolumn{4}{p{0.6\linewidth}}{\cellcolor[rgb]{0.9,0.9,0.9}{%
	\makecell[{{p{\linewidth}}}]{% 
	  \tt {\tiny [A$\langle$GM]}  
	 Penalità: (4,5) non è vuoto, ma contiene 'P'.\\ \tt \\ \tt Avete accumulato collettivamente 6 di 8 penalità.\\ \tt Riprova! 
	  } 
	   } 
	   } 
	 & \\ 
 

    \theutterance \stepcounter{utterance}  

    & \multicolumn{4}{p{0.6\linewidth}}{\cellcolor[rgb]{0.8,1,0.9}{%
	\makecell[{{p{\linewidth}}}]{% 
	  \tt {\tiny [A$\rangle$GM]}  
	 SPOSTA: P, (4,6) 
	  } 
	   } 
	   } 
	 & & \\ 
 

    \theutterance \stepcounter{utterance}  

    & & & \multicolumn{2}{p{0.3\linewidth}}{\cellcolor[rgb]{0.95,0.95,0.95}{%
	\makecell[{{p{\linewidth}}}]{% 
	  \tt {\tiny [GM$|$GM]}  
	 SPOSTA: P, (4,6) 
	  } 
	   } 
	   } 
	 & & \\ 
 

    \theutterance \stepcounter{utterance}  

    & & & \multicolumn{2}{p{0.3\linewidth}}{\cellcolor[rgb]{0.95,0.95,0.95}{%
	\makecell[{{p{\linewidth}}}]{% 
	  \tt {\tiny [GM$|$GM]}  
	 Spostato 'P' in (4,6) con successo. La tua griglia aggiornata appare così:\\ \tt ```\\ \tt  1234567\\ \tt ╔═╤═╤═╤═╗\\ \tt ║L│◌│◌│◌║ 1\\ \tt ║◌└─┼─┴─╢ 2\\ \tt ║◌◌◌│C◌◌║ 3\\ \tt ╟──┬┴───╢ 4\\ \tt ║◌◌│◌◌◌◌║ 5\\ \tt ║◌◌│P◌◌◌║ 6\\ \tt ║◌◌│◌◌◌◌║ 7\\ \tt ╚══╧════╝\\ \tt ```\\ \tt  
	  } 
	   } 
	   } 
	 & & \\ 
 

    \theutterance \stepcounter{utterance}  

    & & \multicolumn{4}{p{0.6\linewidth}}{\cellcolor[rgb]{0.9,0.9,0.9}{%
	\makecell[{{p{\linewidth}}}]{% 
	  \tt {\tiny [GM$\rangle$B]}  
	 Il tuo messaggio è stato inoltrato all'altro giocatore.\\ \tt \\ \tt \\ \tt Stai attualmente giocando il round 5 di massimo 12.\\ \tt Avete accumulato collettivamente 6 di 8 penalità.\\ \tt L'altro giocatore ha spostato un oggetto sulla sua griglia.\\ \tt \\ \tt \\ \tt Qual è il tuo prossimo comando? 
	  } 
	   } 
	   } 
	 & \\ 
 

    \theutterance \stepcounter{utterance}  

    & & & \multicolumn{4}{p{0.6\linewidth}}{\cellcolor[rgb]{1,0.85,0.72}{%
	\makecell[{{p{\linewidth}}}]{% 
	  \tt {\tiny [GM$\langle$B]}  
	 COMUNICA: Grazie per il movimento! Conferma se l'oggetto spostato è in una posizione che si avvicina al nostro obiettivo (C: (4,3), L: (2,1), P: (4,5)). Se non è così, possiamo modificare le coordinate per ottimizzare il piano. 😊 
	  } 
	   } 
	   } 
	  \\ 
 

    \theutterance \stepcounter{utterance}  

    & & & \multicolumn{2}{p{0.3\linewidth}}{\cellcolor[rgb]{0.95,0.95,0.95}{%
	\makecell[{{p{\linewidth}}}]{% 
	  \tt {\tiny [GM$|$GM]}  
	 COMUNICA: Grazie per il movimento! Conferma se l'oggetto spostato è in una posizione che si avvicina al nostro obiettivo (C: (4,3), L: (2,1), P: (4,5)). Se non è così, possiamo modificare le coordinate per ottimizzare il piano. 😊 
	  } 
	   } 
	   } 
	 & & \\ 
 

    \theutterance \stepcounter{utterance}  

    & & \multicolumn{4}{p{0.6\linewidth}}{\cellcolor[rgb]{0.9,0.9,0.9}{%
	\makecell[{{p{\linewidth}}}]{% 
	  \tt {\tiny [A$\langle$GM]}  
	 Spostato 'P' in (4,6) con successo. La tua griglia aggiornata appare così:\\ \tt ```\\ \tt  1234567\\ \tt ╔═╤═╤═╤═╗\\ \tt ║L│◌│◌│◌║ 1\\ \tt ║◌└─┼─┴─╢ 2\\ \tt ║◌◌◌│C◌◌║ 3\\ \tt ╟──┬┴───╢ 4\\ \tt ║◌◌│◌◌◌◌║ 5\\ \tt ║◌◌│P◌◌◌║ 6\\ \tt ║◌◌│◌◌◌◌║ 7\\ \tt ╚══╧════╝\\ \tt ```\\ \tt Stai attualmente giocando il round 5 di massimo 12.\\ \tt Avete accumulato collettivamente 6 di 8 penalità.\\ \tt L'altro giocatore ha inviato questo messaggio:\\ \tt \\ \tt \\ \tt "Grazie per il movimento! Conferma se l'oggetto spostato è in una posizione che si avvicina al nostro obiettivo (C: (4,3), L: (2,1), P: (4,5)). Se non è così, possiamo modificare le coordinate per ottimizzare il piano. 😊"\\ \tt \\ \tt \\ \tt Qual è il tuo prossimo comando? 
	  } 
	   } 
	   } 
	 & \\ 
 

    \theutterance \stepcounter{utterance}  

    & \multicolumn{4}{p{0.6\linewidth}}{\cellcolor[rgb]{0.8,1,0.9}{%
	\makecell[{{p{\linewidth}}}]{% 
	  \tt {\tiny [A$\rangle$GM]}  
	 SPOSTA: P, (4,5) 
	  } 
	   } 
	   } 
	 & & \\ 
 

    \theutterance \stepcounter{utterance}  

    & & & \multicolumn{2}{p{0.3\linewidth}}{\cellcolor[rgb]{0.95,0.95,0.95}{%
	\makecell[{{p{\linewidth}}}]{% 
	  \tt {\tiny [GM$|$GM]}  
	 SPOSTA: P, (4,5) 
	  } 
	   } 
	   } 
	 & & \\ 
 

    \theutterance \stepcounter{utterance}  

    & & & \multicolumn{2}{p{0.3\linewidth}}{\cellcolor[rgb]{0.95,0.95,0.95}{%
	\makecell[{{p{\linewidth}}}]{% 
	  \tt {\tiny [GM$|$GM]}  
	 Spostato 'P' in (4,5) con successo. La tua griglia aggiornata appare così:\\ \tt ```\\ \tt  1234567\\ \tt ╔═╤═╤═╤═╗\\ \tt ║L│◌│◌│◌║ 1\\ \tt ║◌└─┼─┴─╢ 2\\ \tt ║◌◌◌│C◌◌║ 3\\ \tt ╟──┬┴───╢ 4\\ \tt ║◌◌│P◌◌◌║ 5\\ \tt ║◌◌│◌◌◌◌║ 6\\ \tt ║◌◌│◌◌◌◌║ 7\\ \tt ╚══╧════╝\\ \tt ```\\ \tt  
	  } 
	   } 
	   } 
	 & & \\ 
 

    \theutterance \stepcounter{utterance}  

    & & \multicolumn{4}{p{0.6\linewidth}}{\cellcolor[rgb]{0.9,0.9,0.9}{%
	\makecell[{{p{\linewidth}}}]{% 
	  \tt {\tiny [GM$\rangle$B]}  
	 Il tuo messaggio è stato inoltrato all'altro giocatore.\\ \tt \\ \tt \\ \tt Stai attualmente giocando il round 6 di massimo 12.\\ \tt Avete accumulato collettivamente 6 di 8 penalità.\\ \tt L'altro giocatore ha spostato un oggetto sulla sua griglia.\\ \tt \\ \tt \\ \tt Qual è il tuo prossimo comando? 
	  } 
	   } 
	   } 
	 & \\ 
 

    \theutterance \stepcounter{utterance}  

    & & & \multicolumn{4}{p{0.6\linewidth}}{\cellcolor[rgb]{1,0.85,0.72}{%
	\makecell[{{p{\linewidth}}}]{% 
	  \tt {\tiny [GM$\langle$B]}  
	 COMUNICA: Grazie per il movimento! Conferma se l'oggetto spostato è in una posizione che si avvicina al nostro obiettivo (C: (4,3), L: (2,1), P: (4,5)). Se non è così, possiamo modificare le coordinate per ottimizzare il piano. 😊 
	  } 
	   } 
	   } 
	  \\ 
 

    \theutterance \stepcounter{utterance}  

    & & & \multicolumn{2}{p{0.3\linewidth}}{\cellcolor[rgb]{0.95,0.95,0.95}{%
	\makecell[{{p{\linewidth}}}]{% 
	  \tt {\tiny [GM$|$GM]}  
	 COMUNICA: Grazie per il movimento! Conferma se l'oggetto spostato è in una posizione che si avvicina al nostro obiettivo (C: (4,3), L: (2,1), P: (4,5)). Se non è così, possiamo modificare le coordinate per ottimizzare il piano. 😊 
	  } 
	   } 
	   } 
	 & & \\ 
 

    \theutterance \stepcounter{utterance}  

    & & \multicolumn{4}{p{0.6\linewidth}}{\cellcolor[rgb]{0.9,0.9,0.9}{%
	\makecell[{{p{\linewidth}}}]{% 
	  \tt {\tiny [A$\langle$GM]}  
	 Spostato 'P' in (4,5) con successo. La tua griglia aggiornata appare così:\\ \tt ```\\ \tt  1234567\\ \tt ╔═╤═╤═╤═╗\\ \tt ║L│◌│◌│◌║ 1\\ \tt ║◌└─┼─┴─╢ 2\\ \tt ║◌◌◌│C◌◌║ 3\\ \tt ╟──┬┴───╢ 4\\ \tt ║◌◌│P◌◌◌║ 5\\ \tt ║◌◌│◌◌◌◌║ 6\\ \tt ║◌◌│◌◌◌◌║ 7\\ \tt ╚══╧════╝\\ \tt ```\\ \tt Stai attualmente giocando il round 6 di massimo 12.\\ \tt Avete accumulato collettivamente 6 di 8 penalità.\\ \tt L'altro giocatore ha inviato questo messaggio:\\ \tt \\ \tt \\ \tt "Grazie per il movimento! Conferma se l'oggetto spostato è in una posizione che si avvicina al nostro obiettivo (C: (4,3), L: (2,1), P: (4,5)). Se non è così, possiamo modificare le coordinate per ottimizzare il piano. 😊"\\ \tt \\ \tt \\ \tt Qual è il tuo prossimo comando? 
	  } 
	   } 
	   } 
	 & \\ 
 

    \theutterance \stepcounter{utterance}  

    & \multicolumn{4}{p{0.6\linewidth}}{\cellcolor[rgb]{0.8,1,0.9}{%
	\makecell[{{p{\linewidth}}}]{% 
	  \tt {\tiny [A$\rangle$GM]}  
	 SPOSTA: C, (4,3) 
	  } 
	   } 
	   } 
	 & & \\ 
 

    \theutterance \stepcounter{utterance}  

    & & & \multicolumn{2}{p{0.3\linewidth}}{\cellcolor[rgb]{0.95,0.95,0.95}{%
	\makecell[{{p{\linewidth}}}]{% 
	  \tt {\tiny [GM$|$GM]}  
	 SPOSTA: C, (4,3) 
	  } 
	   } 
	   } 
	 & & \\ 
 

    \theutterance \stepcounter{utterance}  

    & & & \multicolumn{2}{p{0.3\linewidth}}{\cellcolor[rgb]{0.95,0.95,0.95}{%
	\makecell[{{p{\linewidth}}}]{% 
	  \tt {\tiny [GM$|$GM]}  
	 Penalità: (4,3) non è vuoto, ma contiene '│'.\\ \tt \\ \tt Avete accumulato collettivamente 7 di 8 penalità.\\ \tt Riprova! 
	  } 
	   } 
	   } 
	 & & \\ 
 

    \theutterance \stepcounter{utterance}  

    & & \multicolumn{4}{p{0.6\linewidth}}{\cellcolor[rgb]{0.9,0.9,0.9}{%
	\makecell[{{p{\linewidth}}}]{% 
	  \tt {\tiny [A$\langle$GM]}  
	 Penalità: (4,3) non è vuoto, ma contiene '│'.\\ \tt \\ \tt Avete accumulato collettivamente 7 di 8 penalità.\\ \tt Riprova! 
	  } 
	   } 
	   } 
	 & \\ 
 

    \theutterance \stepcounter{utterance}  

    & \multicolumn{4}{p{0.6\linewidth}}{\cellcolor[rgb]{0.8,1,0.9}{%
	\makecell[{{p{\linewidth}}}]{% 
	  \tt {\tiny [A$\rangle$GM]}  
	 SPOSTA: C, (4,2) 
	  } 
	   } 
	   } 
	 & & \\ 
 

    \theutterance \stepcounter{utterance}  

    & & & \multicolumn{2}{p{0.3\linewidth}}{\cellcolor[rgb]{0.95,0.95,0.95}{%
	\makecell[{{p{\linewidth}}}]{% 
	  \tt {\tiny [GM$|$GM]}  
	 SPOSTA: C, (4,2) 
	  } 
	   } 
	   } 
	 & & \\ 
 

    \theutterance \stepcounter{utterance}  

    & & & \multicolumn{2}{p{0.3\linewidth}}{\cellcolor[rgb]{0.95,0.95,0.95}{%
	\makecell[{{p{\linewidth}}}]{% 
	  \tt {\tiny [GM$|$GM]}  
	 Penalità: (4,2) non è vuoto, ma contiene '┼'.\\ \tt \\ \tt Avete accumulato collettivamente 8 di 8 penalità.\\ \tt Riprova! 
	  } 
	   } 
	   } 
	 & & \\ 
 

    \theutterance \stepcounter{utterance}  

    & & \multicolumn{4}{p{0.6\linewidth}}{\cellcolor[rgb]{0.9,0.9,0.9}{%
	\makecell[{{p{\linewidth}}}]{% 
	  \tt {\tiny [A$\langle$GM]}  
	 Penalità: (4,2) non è vuoto, ma contiene '┼'.\\ \tt \\ \tt Avete accumulato collettivamente 8 di 8 penalità.\\ \tt Riprova! 
	  } 
	   } 
	   } 
	 & \\ 
 

    \theutterance \stepcounter{utterance}  

    & \multicolumn{4}{p{0.6\linewidth}}{\cellcolor[rgb]{0.8,1,0.9}{%
	\makecell[{{p{\linewidth}}}]{% 
	  \tt {\tiny [A$\rangle$GM]}  
	 COMUNICA: Finito? No, non abbiamo ancora raggiunto l'obiettivo. Le penalità sono esaurite, ma possiamo ancora provare a correggere gli errori. Suggerimento: verifica le posizioni degli oggetti e cerca percorsi alternativi per spostarli senza collisioni. 😊 
	  } 
	   } 
	   } 
	 & & \\ 
 

    \theutterance \stepcounter{utterance}  

    & & & \multicolumn{2}{p{0.3\linewidth}}{\cellcolor[rgb]{0.95,0.95,0.95}{%
	\makecell[{{p{\linewidth}}}]{% 
	  \tt {\tiny [GM$|$GM]}  
	 COMUNICA: Finito? No, non abbiamo ancora raggiunto l'obiettivo. Le penalità sono esaurite, ma possiamo ancora provare a correggere gli errori. Suggerimento: verifica le posizioni degli oggetti e cerca percorsi alternativi per spostarli senza collisioni. 😊 
	  } 
	   } 
	   } 
	 & & \\ 
 

    \theutterance \stepcounter{utterance}  

    & & \multicolumn{4}{p{0.6\linewidth}}{\cellcolor[rgb]{0.9,0.9,0.9}{%
	\makecell[{{p{\linewidth}}}]{% 
	  \tt {\tiny [GM$\rangle$B]}  
	 Il tuo messaggio è stato inoltrato all'altro giocatore.\\ \tt \\ \tt \\ \tt Stai attualmente giocando il round 7 di massimo 12.\\ \tt Avete accumulato collettivamente 8 di 8 penalità.\\ \tt L'altro giocatore ha inviato questo messaggio:\\ \tt \\ \tt \\ \tt "Finito? No, non abbiamo ancora raggiunto l'obiettivo. Le penalità sono esaurite, ma possiamo ancora provare a correggere gli errori. Suggerimento: verifica le posizioni degli oggetti e cerca percorsi alternativi per spostarli senza collisioni. 😊"\\ \tt \\ \tt \\ \tt Qual è il tuo prossimo comando? 
	  } 
	   } 
	   } 
	 & \\ 
 

    \theutterance \stepcounter{utterance}  

    & & & \multicolumn{4}{p{0.6\linewidth}}{\cellcolor[rgb]{1,0.85,0.72}{%
	\makecell[{{p{\linewidth}}}]{% 
	  \tt {\tiny [GM$\langle$B]}  
	 COMUNICA: Accordo sul suggerimento! Verificherò le posizioni degli oggetti sul mio tavolo e cercherò percorsi alternativi per spostare 'C', 'L' e 'P' verso (4,3), (2,1) e (4,5) rispettivamente, evitando collisioni. Proponiamo di iniziare con il movimento di 'L' verso (2,1) se è fattibile. 😊 
	  } 
	   } 
	   } 
	  \\ 
 

    \theutterance \stepcounter{utterance}  

    & & & \multicolumn{2}{p{0.3\linewidth}}{\cellcolor[rgb]{0.95,0.95,0.95}{%
	\makecell[{{p{\linewidth}}}]{% 
	  \tt {\tiny [GM$|$GM]}  
	 COMUNICA: Accordo sul suggerimento! Verificherò le posizioni degli oggetti sul mio tavolo e cercherò percorsi alternativi per spostare 'C', 'L' e 'P' verso (4,3), (2,1) e (4,5) rispettivamente, evitando collisioni. Proponiamo di iniziare con il movimento di 'L' verso (2,1) se è fattibile. 😊 
	  } 
	   } 
	   } 
	 & & \\ 
 

    \theutterance \stepcounter{utterance}  

    & & \multicolumn{4}{p{0.6\linewidth}}{\cellcolor[rgb]{0.9,0.9,0.9}{%
	\makecell[{{p{\linewidth}}}]{% 
	  \tt {\tiny [A$\langle$GM]}  
	 Il tuo messaggio è stato inoltrato all'altro giocatore.\\ \tt \\ \tt \\ \tt Stai attualmente giocando il round 7 di massimo 12.\\ \tt Avete accumulato collettivamente 8 di 8 penalità.\\ \tt L'altro giocatore ha inviato questo messaggio:\\ \tt \\ \tt \\ \tt "Accordo sul suggerimento! Verificherò le posizioni degli oggetti sul mio tavolo e cercherò percorsi alternativi per spostare 'C', 'L' e 'P' verso (4,3), (2,1) e (4,5) rispettivamente, evitando collisioni. Proponiamo di iniziare con il movimento di 'L' verso (2,1) se è fattibile. 😊"\\ \tt \\ \tt \\ \tt Qual è il tuo prossimo comando? 
	  } 
	   } 
	   } 
	 & \\ 
 

    \theutterance \stepcounter{utterance}  

    & \multicolumn{4}{p{0.6\linewidth}}{\cellcolor[rgb]{0.8,1,0.9}{%
	\makecell[{{p{\linewidth}}}]{% 
	  \tt {\tiny [A$\rangle$GM]}  
	 SPOSTA: L, (2,1) 
	  } 
	   } 
	   } 
	 & & \\ 
 

    \theutterance \stepcounter{utterance}  

    & & & \multicolumn{2}{p{0.3\linewidth}}{\cellcolor[rgb]{0.95,0.95,0.95}{%
	\makecell[{{p{\linewidth}}}]{% 
	  \tt {\tiny [GM$|$GM]}  
	 SPOSTA: L, (2,1) 
	  } 
	   } 
	   } 
	 & & \\ 
 

    \theutterance \stepcounter{utterance}  

    & & & \multicolumn{2}{p{0.3\linewidth}}{\cellcolor[rgb]{0.95,0.95,0.95}{%
	\makecell[{{p{\linewidth}}}]{% 
	  \tt {\tiny [GM$|$GM]}  
	 Penalità: (2,1) non è vuoto, ma contiene '│'.\\ \tt \\ \tt Avete accumulato collettivamente 9 di 8 penalità.\\ \tt Riprova! 
	  } 
	   } 
	   } 
	 & & \\ 
 

    \theutterance \stepcounter{utterance}  

    & & & \multicolumn{2}{p{0.3\linewidth}}{\cellcolor[rgb]{0.95,0.95,0.95}{%
	\makecell[{{p{\linewidth}}}]{% 
	  \tt {\tiny [GM$|$GM]}  
	 Maximum number of penalties exceeded 
	  } 
	   } 
	   } 
	 & & \\ 
 

    \theutterance \stepcounter{utterance}  

    & & & \multicolumn{2}{p{0.3\linewidth}}{\cellcolor[rgb]{0.95,0.95,0.95}{%
	\makecell[{{p{\linewidth}}}]{% 
	  \tt {\tiny [GM$|$GM]}  
	 * success: False\\ \tt * lose: True\\ \tt * aborted: True\\ \tt -------\\ \tt * Shifts: 2.00\\ \tt * Max Shifts: 4.00\\ \tt * Min Shifts: 2.00\\ \tt * End Distance Sum: 7.41\\ \tt * Init Distance Sum: 13.65\\ \tt * Expected Distance Sum: 12.57\\ \tt * Penalties: 9.00\\ \tt * Max Penalties: 8.00\\ \tt * Rounds: 7.00\\ \tt * Max Rounds: 12.00\\ \tt * Object Count: 3.00\\ \tt  
	  } 
	   } 
	   } 
	 & & \\ 
 

    \theutterance \stepcounter{utterance}  

    & & & \multicolumn{2}{p{0.3\linewidth}}{\cellcolor[rgb]{0.95,0.95,0.95}{%
	\makecell[{{p{\linewidth}}}]{% 
	  \tt {\tiny [GM$|$GM]}  
	 * Main Score: 0.00\\ \tt \\ \tt -------\\ \tt * Distance Score: 0.43\\ \tt * Consistency Score: 1.00\\ \tt * Coverage Score: 0.25\\ \tt * Penalty Score: 0.00\\ \tt * Alternative Penalty Score: 0.00\\ \tt * Alternative Main Score: 0.00\\ \tt \\ \tt -------\\ \tt * Shifts: 2.00\\ \tt * Max Shifts: 4.00\\ \tt * Min Shifts: 2.00\\ \tt * End Distance Sum: 7.41\\ \tt * Init Distance Sum: 13.65\\ \tt * Expected Distance Sum: 12.57\\ \tt * Penalties: 9.00\\ \tt * Max Penalties: 8.00\\ \tt * Rounds: 7.00\\ \tt * Max Rounds: 12.00\\ \tt * Object Count: 3.00\\ \tt  
	  } 
	   } 
	   } 
	 & & \\ 
 

\end{supertabular}
}

\end{document}
